The operating constraints that shaped every experiment are documented
in the charter and decision log \cite{overview,decisionlog}. They are
restated here because they are load-bearing for the interpretation of
every result.

\paragraph{Specified interfaces, not architectures.}
The project defines the sensory front end (input channels with
deterministic Poisson spike trains per presentation), the output
boundary (ordinary plastic neurons whose spikes are additionally
emitted as output activity --- a role tag, not a cell type), the
environment/curriculum (an ordered schedule of pattern presentations),
resource limits (neuron, synapse, per-neuron budget caps), and the
plasticity laws (STDP, weight normalization, candidate/prune machinery,
anti-Hebbian inhibition). It never defines assemblies, readout circuits,
or any intermediate organization.

\paragraph{Local rules.}
Every mechanism is required to read only per-neuron or per-synapse
state. This is enforced loosely in design and strictly at review: the
CLLA falsifier F5 requires static conformance --- consolidation and
allocation decision paths may not read pattern identity, channel-group
membership tests, presentation counts, or any global state
\cite{cllaprot,cllarecord}.

\paragraph{Finite resources.}
A neuron owns a total excitatory budget \(t_e = 0.8\) (weight units, M2
normalization) and a candidate pool of six slots; consolidation may
protect at most \(p_{\max} t_e = 0.60\) of the budget; budgets bind
synapse creation (M5, \(b_e = 40\) excitatory afferents per neuron).
Resource limits are treated as part of the specification, not as an
implementation nuisance.

\paragraph{No semantic supervision.}
All presented patterns are examples; the curriculum is a configurable
schedule, not a set of preinstalled concepts. The organism is never told
which pattern it is seeing, whether a response was correct, or that a
pattern is ``novel.''

\paragraph{Reproducibility and freeze discipline.}
Protocols are committed before implementation; amendments are logged
and user-approved; every execution is byte-deterministic under a seeded
Xoshiro generator; a flag-off ``identity gate'' must reproduce the
committed baseline byte-for-byte (event-stream hash and snapshot
frames) before any flag-on run is interpreted; and negative results are
preserved with their runs \cite{overview,cllaprot}. Several protocol
defects were discovered only by this discipline and are reported as
such (\S\ref{sec:clla}, \S\ref{sec:method}).

\begin{figure}[t]
\centering
\resizebox{\textwidth}{!}{%
% Figure 1 — ANIMA architecture + experimental philosophy (schematic).
% Every annotation in the right panel is established in the cited records.
\begin{tikzpicture}[
  box/.style={draw, rounded corners=2pt, minimum height=6mm, font=\small,
              align=center, inner sep=3pt},
  spec/.style={fill=blue!8}, emerge/.style={fill=green!8},
  arrow/.style={-{Stealth[length=2.2mm]}, thick}
]
  % layer 1: input
  \node[box, spec] (inp) at (0, 0) {24 input channels\\groups A / C / B};
  % layer 2: network
  \node[box, spec] (net) at (4.2, 0) {LIF neurons\\40 internal + 12 output};
  % layer 3: plasticity (box to the right)
  \node[box, spec, minimum width=3.2cm] (plast) at (9.0, 0) {local plasticity\\STDP, M2 budget $t_e$,\\M3/M4 candidates/prune,\\M6 inhibition (per synapse)};
  \draw[arrow] (inp) -- node[above, font=\small] {deterministic Poisson, 20 Hz} (net);
  \draw[arrow] (net) -- node[above, font=\small] {weights $w$} (plast);
  \draw[arrow] (plast) to[bend left=30] node[below, font=\small] {budget} (net);
  % environment / curriculum
  \node[box, spec, minimum width=8.9cm] (env) at (4.6, -1.5) {curriculum (frozen schedule): isolated / blocked (bac, bca) / alternating (il); 2 s cadence};
  \draw[arrow] (env) -- (inp);

  % right panel: philosophy
  \node[box, emerge, font=\small, minimum width=6.9cm, text width=6.5cm] (p1) at (13.4, 1.15)
    {\textbf{Specified:} I/O boundary, environment, curriculum,\\resource caps, local plasticity laws};
  \node[box, emerge, font=\small, minimum width=6.9cm, text width=6.5cm] (p2) at (13.4, -0.35)
    {\textbf{Not specified:} architecture, assemblies, readout,\\task semantics, labels, reward};
  \node[box, font=\small, minimum width=6.9cm, text width=6.5cm] (p3) at (13.4, -1.85)
    {\textbf{Must emerge:} organization $\to$ persistence $\to$\\retrieval (the Phase I ladder)};

  % bottom: causal narrowing strip
  \node[box, minimum width=13.6cm, text width=13.2cm, fill=gray!6] (strip) at (6.8, -3.6)
    {\footnotesize
    organization $\to$ transient representation $\to$ episode-identity loss
    $\to$ intrinsic-\(u\) poverty $\to$ synaptic store $\to$ interference
    $\to$ protection $\to$ stability $\to$ allocation $\to$
    first-exposure allocation $\to$ recruitment \;\; (details \S\ref{sec:allocation}--\ref{sec:rg})};
\end{tikzpicture}
}
\caption{\textbf{ANIMA architecture and experimental philosophy.}
Top: the specified boundary --- 24 input channels (groups A, C, B),
64 LIF neurons (40 internal, 12 output plus input neurons), sparse
seeded wiring, local plasticity laws, deterministic curriculum. Bottom
left: the charter---what is specified vs.\ what must emerge. Bottom
right: the phase progression --- the causal narrowing of the persistent-
learning failure reported in this paper (detail in \S\ref{sec:synthesis}).}
\label{fig:architecture}
\end{figure}